\chapter{Laws of Large Numbers}
\begin{center}\itshape
A casino offers a game: a coin is tossed until it shows heads, and if that takes $k$ tosses you win $2^k$ euros.\\
The expected win is infinite. What is a fair price for playing $n$ times?\\
About $n\log_2n$ euros, and the law of large numbers proves it, in probability but not almost surely.
\end{center}

Jakob Bernoulli died in 1705 with his book on probability unfinished. His nephew Nicolaus published it in 1713 as \emph{Ars
Conjectandi}, and its fourth part contains the first limit theorem of probability. Draw with replacement from an urn in which
the white pebbles make up three fifths of the total. How many draws are needed so that the observed proportion of white is
within $\frac1{50}$ of $\frac35$, with odds of a thousand to one? Bernoulli proved that a finite number suffices, and computed
one: $25\,550$. The number is absurdly large, because his inequalities were crude, but the principle was new: frequencies converge
to probabilities, and one can say how fast. Simeon-Denis Poisson named the principle \emph{la loi des grands nombres} in 1837.

For two centuries ``converge'' meant ``in probability''. Borel's paper of 1909 (Chapter~1) asked for more: that the
frequency of heads in one infinite sequence of tosses tends to $\frac12$, for almost every sequence. This is the \emph{strong}
law, and it is genuinely stronger. Francesco Cantelli proved it in 1917 for variables with four moments, the proof in Tome~3.
Aleksandr Khintchine proved the weak law for i.i.d.\ variables with a finite mean in 1929, and Kolmogorov the strong law in the same
generality in 1930, with a proof through random series that we shall meet in Section~\ref{sec:threeseries}. Fifty years later, in 1981,
Nasrollah Etemadi published a proof in four pages that needs only pairwise independence, and it is the one this chapter gives.

Then there is the casino of the epigraph, the \emph{St Petersburg game}. Nicolaus Bernoulli, the same nephew, proposed it in a letter
to Pierre R\'emond de Montmort in 1713; his cousin Daniel Bernoulli analysed it in the proceedings of the St Petersburg academy in 1738
and gave it its name. The expected gain $\sum_k2^k2^{-k}$ is infinite, yet nobody would pay much to play. Daniel's answer was the
utility of money, which founded economics' theory of choice. William Feller's answer, in 1945, was a law of large numbers: the
total gain after $n$ games is $n\log_2n$ to first order, in probability. Fair entrance fees exist after all; they just grow with $n$.
And the almost sure version fails: infinitely often, the player is far ahead (Exercise~\ref{ex:stpeteas}).

The chapter has three tools. \emph{Truncation}, which cuts off the rare large values that make the variance infinite. Kolmogorov's
\emph{three-series theorem}, which says exactly when a series of independent terms converges. And \emph{monotone interpolation}, which
passes from a sparse geometric subsequence of times to all times. With them we prove Feller's weak law and the St Petersburg price,
Etemadi's strong law, the three-series theorem with its converse, and the strong law of renewal theory.

\begin{reminder}[from Tome 3 and Chapter 1]
\emph{Chebyshev's inequality} and the weak law for uncorrelated variables with bounded variances (Tome~3, Theorem~4.2): for pairwise independent
variables, $\operatorname{Var}(X_1+\dots+X_n)=\sum\operatorname{Var}X_k$. \emph{Optional stopping} (Tome~3, Theorem~10.12): if $(M_k)$ is a martingale and
$\tau\le n$ a stopping time, $\mathbb E M_\tau=\mathbb E M_0$. \emph{Borel--Cantelli} (Tome~3, Theorem~4.11). \emph{Chapter~1}: the zero--one law,
Kolmogorov's $L^2$ criterion (Proposition~\ref{prop:l2series}), and the fact that a convergent series of independent terms with
$\mathbb E|X_1|=\infty$ cannot satisfy $S_n/n\to c$ (Exercise~\ref{ex:infmean}). For $Y\ge0$, $\sum_{n\ge1}P(Y>n)\le\mathbb E Y\le\sum_{n\ge0}P(Y>n)$.
\end{reminder}

\section{Truncation and the weak law}

When the variance is infinite, Chebyshev's inequality is useless as it stands. The remedy is to cut: replace $X$ by $X\mathbf 1_{|X|\le b}$,
which has all moments, and control separately the probability that the cut made a difference. The following form of the idea works for
triangular arrays, where the summands may change with $n$.

\begin{theorem}[weak law for triangular arrays]\label{thm:wlltriangular}
For each $n$, let $X_{n,1},\dots,X_{n,n}$ be independent, let $b_n>0$, and let $Y_{n,k}=X_{n,k}\mathbf 1_{|X_{n,k}|\le b_n}$. Suppose
\[
\text{(a)}\ \sum_{k=1}^nP(|X_{n,k}|>b_n)\to0,\qquad\text{(b)}\ \frac1{b_n^2}\sum_{k=1}^n\mathbb E\,Y_{n,k}^2\to0.
\]
Then, with $S_n=\sum_kX_{n,k}$ and $a_n=\sum_k\mathbb E\,Y_{n,k}$, we have $(S_n-a_n)/b_n\to0$ in probability.
\end{theorem}
\begin{proof}
Let $T_n=\sum_kY_{n,k}$. Then $P(S_n\ne T_n)\le\sum_kP(|X_{n,k}|>b_n)\to0$ by~(a). By Chebyshev and independence,
\[
P\Bigl(\Bigl|\frac{T_n-a_n}{b_n}\Bigr|>\varepsilon\Bigr)\le\frac{\operatorname{Var}T_n}{\varepsilon^2b_n^2}\le\frac1{\varepsilon^2b_n^2}\sum_k\mathbb E\,Y_{n,k}^2\to0
\]
by~(b). So $P(|S_n-a_n|>\varepsilon b_n)\le P(S_n\ne T_n)+P(|T_n-a_n|>\varepsilon b_n)\to0$.
\end{proof}

For i.i.d.\ variables with a finite mean, take $b_n=n$: the sums in (a) and (b) are $nP(|X|>n)$ and $\frac1n\mathbb E[X^2\mathbf 1_{|X|\le n}]$, both
of which tend to $0$ (Exercise~\ref{ex:feller}), and $a_n/n=\mathbb E[X\mathbf 1_{|X|\le n}]\to\mathbb E X$: this is Khintchine's weak law. The
St Petersburg game has no mean, but the theorem still applies, with a different scale.

\begin{theorem}[Feller, 1945]\label{thm:stpetersburg}
Let $X_1,X_2,\dots$ be i.i.d.\ with $P(X=2^j)=2^{-j}$ for $j\ge1$, and $S_n=X_1+\dots+X_n$. Then
\[
\frac{S_n}{n\log_2n}\to1\quad\text{in probability.}
\]
\end{theorem}
\begin{proof}
Apply Theorem~\ref{thm:wlltriangular} with $X_{n,k}=X_k$ and $b_n=n\log_2n$. Since $P(X\ge2^j)=2^{1-j}$, we have $P(X>x)\le2/x$ for all
$x>0$, so (a) is $nP(X>b_n)\le2n/b_n=2/\log_2n\to0$. For (b),
\[
\mathbb E[X^2\mathbf 1_{X\le b_n}]=\sum_{j:\,2^j\le b_n}2^{2j}2^{-j}=\sum_{j\le\log_2b_n}2^j\le2b_n,
\]
so $nb_n^{-2}\mathbb E[X^2\mathbf 1_{X\le b_n}]\le2n/b_n\to0$. Finally $a_n=n\,\mathbb E[X\mathbf 1_{X\le b_n}]=n\sum_{j\le\log_2b_n}1=n\lfloor\log_2b_n\rfloor$,
and $\log_2b_n=\log_2n+\log_2\log_2n$, so $a_n/b_n\to1$. Hence $S_n/b_n-a_n/b_n\to0$ in probability, and $S_n/b_n\to1$.
\end{proof}

\begin{figure}[ht]
\centering\fig{ch02-petersburg}
\caption{Five players of the St Petersburg game: $S_n/(n\log_2n)$ against $n$ on a logarithmic scale, clipped at $4$. The ratio drifts towards $1$
(Theorem~\ref{thm:stpetersburg}) but every so often a single huge prize throws it upwards; one of these five paths reaches $7$ after $n=1000$.
The spikes never stop (Exercise~\ref{ex:stpeteas}).}
\label{fig:petersburg}
\end{figure}

\begin{think}
In the St Petersburg game, how long must you play before a single prize exceeds everything won so far? If the largest of $n$ prizes is about $n$
(why?), and the total is about $n\log_2n$, what fraction of your total comes from your best game? Now explain the jumps in Figure~
\ref{fig:petersburg}.
\end{think}

\section{The strong law}

\begin{theorem}[Etemadi, 1981]\label{thm:etemadi}
Let $X_1,X_2,\dots$ be pairwise independent and identically distributed, with $\mathbb E|X_1|<\infty$ and $\mu=\mathbb E X_1$. Then
$S_n/n\to\mu$ almost surely.
\end{theorem}
\begin{proof}
\emph{Reduction.} The positive parts $X_k^+$ are again pairwise independent and identically distributed, and so are the negative parts, and
$S_n=\sum X_k^+-\sum X_k^-$. So we may assume $X_k\ge0$.

\emph{Truncation.} Let $Y_k=X_k\mathbf 1_{X_k\le k}$ and $T_n=Y_1+\dots+Y_n$. Since $\sum_kP(X_k\ne Y_k)=\sum_{k\ge1}P(X_1>k)\le\mathbb E X_1<\infty$,
Borel--Cantelli says that $X_k=Y_k$ for all large $k$, almost surely. Then $S_n-T_n$ is eventually constant, so it suffices to show
$T_n/n\to\mu$. Also $\mathbb E Y_k=\mathbb E[X_1\mathbf 1_{X_1\le k}]\to\mu$ by monotone convergence, so $\mathbb E T_n/n\to\mu$ (Ces\`aro).

\emph{A variance bound.} The $Y_k$ are pairwise independent, so $\operatorname{Var}T_n=\sum_{k\le n}\operatorname{Var}Y_k\le\sum_{k\le n}\mathbb E Y_k^2$.
We shall need
\begin{equation}\label{eq:etemadisum}
\sum_{k\ge1}\frac{\mathbb E Y_k^2}{k^2}\le2\,\mathbb E X_1.
\end{equation}
Indeed $\sum_kk^{-2}\mathbb E Y_k^2=\mathbb E\bigl[X_1^2\sum_{k\ge\max(X_1,1)}k^{-2}\bigr]$, and $\sum_{k\ge y}k^{-2}\le\frac1{y^2}+\int_y^\infty
\frac{\dd x}{x^2}\le\frac2y$ for $y\ge1$; so the bracket is at most $2X_1^2/\max(X_1,1)\le2X_1$.

\emph{A geometric subsequence.} Fix $\alpha>1$ and let $k_n=\lfloor\alpha^n\rfloor$. For $\varepsilon>0$, Chebyshev gives
\[
\sum_nP\bigl(|T_{k_n}-\mathbb E T_{k_n}|>\varepsilon k_n\bigr)\le\frac1{\varepsilon^2}\sum_n\frac1{k_n^2}\sum_{m\le k_n}\mathbb E Y_m^2=
\frac1{\varepsilon^2}\sum_m\mathbb E Y_m^2\sum_{n:\,k_n\ge m}\frac1{k_n^2}.
\]
Since $k_n\ge\alpha^n/2$, the inner sum is at most $4\sum_{n:\,\alpha^n\ge m}\alpha^{-2n}\le\frac{4}{(1-\alpha^{-2})m^2}$, and by~\eqref{eq:etemadisum}
the whole is finite. By Borel--Cantelli, $(T_{k_n}-\mathbb E T_{k_n})/k_n\to0$ almost surely, and so $T_{k_n}/k_n\to\mu$ almost surely.

\emph{Interpolation.} Here positivity is used. For $k_n\le m\le k_{n+1}$, since $T$ is non-decreasing,
\[
\frac{T_{k_n}}{k_{n+1}}\le\frac{T_m}m\le\frac{T_{k_{n+1}}}{k_n}.
\]
As $k_{n+1}/k_n\to\alpha$, we get $\mu/\alpha\le\liminf T_m/m\le\limsup T_m/m\le\alpha\mu$ almost surely. Letting $\alpha\downarrow1$ along a
sequence, $T_m/m\to\mu$ almost surely.
\end{proof}

The proof never uses more than pairs of variables at a time: Chebyshev's inequality only sees variances. Everything else is Borel--Cantelli,
which needs no independence in its first half. The integrability assumption is necessary, by Exercise~\ref{ex:infmean} of Chapter~1: if
$\mathbb E|X_1|=\infty$, then $\limsup|S_n|/n=\infty$. For non-negative variables, $S_n/n\to\infty$ (Exercise~\ref{ex:infmeanpos}).

\begin{figure}[ht]
\centering\fig{ch02-pareto}
\caption{Running averages of five samples from the Pareto law $P(X>x)=x^{-3/2}$, $x\ge1$, which has mean $3$ and infinite variance. Each path
converges to $3$ (Theorem~\ref{thm:etemadi}), but slowly and with upward jumps: with infinite variance there is no $\frac1{\sqrt n}$ rate.
Chapter~5 finds the true scale, $n^{-1/3}$, and the stable law that replaces the Gaussian.}
\label{fig:pareto}
\end{figure}

\section{Kolmogorov's three-series theorem}\label{sec:threeseries}

Chapter~1 settled random signs. For general independent terms, Kolmogorov found three conditions that together are necessary and sufficient
for convergence. The converse direction needs a lower bound for the maximum of a walk with bounded steps.

\begin{lemma}[Kolmogorov's other inequality]\label{lem:kolmogorovlower}
Let $X_1,\dots,X_n$ be independent, with $\mathbb E X_k=0$ and $|X_k|\le c$. Let $S_k=X_1+\dots+X_k$ and $s_n^2=\operatorname{Var}S_n$. Then for
$a>0$,
\[
P\Bigl(\max_{k\le n}|S_k|\le a\Bigr)\le\frac{(a+c)^2}{s_n^2}.
\]
\end{lemma}
\begin{proof}
Let $s_k^2=\operatorname{Var}S_k$. The process $M_k=S_k^2-s_k^2$ is a martingale: $\mathbb E[S_{k+1}^2-S_k^2\mid\mathcal F_k]=\mathbb E[2S_kX_{k+1}+X_{k+1}^2
\mid\mathcal F_k]=\operatorname{Var}X_{k+1}$. Let $\tau$ be the first $k\le n$ with $|S_k|>a$, and $\tau=n$ if there is none; it is a bounded
stopping time, so optional stopping gives $\mathbb E S_\tau^2=\mathbb E s_\tau^2$. Now $|S_\tau|\le a+c$: either $\tau$ is the first exit, and then
$|S_{\tau-1}|\le a$ and $|X_\tau|\le c$, or there is no exit and $|S_n|\le a$. So $\mathbb E s_\tau^2\le(a+c)^2$. On the event
$A=\{\max_{k\le n}|S_k|\le a\}$ we have $\tau=n$, so $\mathbb E s_\tau^2\ge s_n^2P(A)$.
\end{proof}

\begin{theorem}[three-series theorem, Kolmogorov]\label{thm:threeseries}
Let $X_1,X_2,\dots$ be independent, $c>0$, and $Y_n=X_n\mathbf 1_{|X_n|\le c}$. Then $\sum X_n$ converges almost surely if and only if the three
series
\[
\text{(i)}\ \sum_nP(|X_n|>c),\qquad\text{(ii)}\ \sum_n\mathbb E Y_n,\qquad\text{(iii)}\ \sum_n\operatorname{Var}Y_n
\]
all converge. Otherwise it diverges almost surely.
\end{theorem}
\begin{proof}
\emph{Sufficiency.} By (iii) and Proposition~\ref{prop:l2series}, $\sum(Y_n-\mathbb E Y_n)$ converges almost surely; with (ii), so does $\sum Y_n$.
By (i) and Borel--Cantelli, $X_n=Y_n$ for all large $n$, almost surely, so $\sum X_n$ converges.

\emph{Necessity.} Convergence is a tail event, so if it does not hold almost surely it fails almost surely. Suppose $\sum X_n$ converges almost surely.
Then $X_n\to0$, so $\{|X_n|>c\}$ happens only finitely often; these events are independent, so by the second Borel--Cantelli lemma (i) holds. As above,
$\sum Y_n$ then converges almost surely.

To reach the variances, symmetrise: on a product space, let $(Y'_n)$ be an independent copy of $(Y_n)$ (Tome~6, Theorem~6.17), and let
$Z_n=Y_n-Y'_n$. The $Z_n$ are independent, centred, $|Z_n|\le2c$, $\operatorname{Var}Z_n=2\operatorname{Var}Y_n$, and $\sum Z_n$ converges almost
surely. Its partial sums $U_n$ are therefore bounded almost surely, so there is $a$ with $P(\sup_n|U_n|\le a)\ge\frac12$. By
Lemma~\ref{lem:kolmogorovlower}, for every $n$,
\[
2\sum_{k\le n}\operatorname{Var}Y_k=\operatorname{Var}U_n\le\frac{(a+2c)^2}{P(\max_{k\le n}|U_k|\le a)}\le2(a+2c)^2.
\]
This is (iii). Then $\sum(Y_n-\mathbb E Y_n)$ converges almost surely by Proposition~\ref{prop:l2series}, and subtracting it from the convergent
$\sum Y_n$ shows that the numerical series $\sum\mathbb E Y_n$ converges: (ii).
\end{proof}

The theorem holds for every $c>0$: if the three series converge for one $c$ they converge for all. It also gives a complete answer to Chapter~1's
question about random signs, now for any law: for instance $\sum X_n/n$ converges almost surely for i.i.d.\ centred $X_n$ with $\mathbb E[|X_1|\log^+|X_1|]<\infty$
(Exercise~\ref{ex:seriesiid}), even when the variance is infinite, while a finite mean alone does not suffice. And it gives Kolmogorov's own route to the strong law, through the following lemma.

\begin{lemma}[Kronecker]\label{lem:kronecker}
If $0<b_n\uparrow\infty$ and $\sum x_n/b_n$ converges, then $\frac1{b_n}\sum_{k\le n}x_k\to0$.
\end{lemma}
\begin{proof}
Let $r_n=\sum_{k>n}x_k/b_k$, so $r_n\to0$ and $x_k=b_k(r_{k-1}-r_k)$. Summation by parts gives
\[
\frac1{b_n}\sum_{k=1}^nx_k=\frac1{b_n}\Bigl(b_1r_0-b_nr_n+\sum_{k=1}^{n-1}(b_{k+1}-b_k)r_k\Bigr).
\]
The first term tends to $0$ and so does the second. The last is a weighted average of the $r_k$ with weights $(b_{k+1}-b_k)/b_n$, whose total is
$(b_n-b_1)/b_n\le1$ and which put vanishing weight on each fixed $k$; since $r_k\to0$, it tends to $0$.
\end{proof}

\begin{remark}[Kolmogorov's proof of the strong law]
For i.i.d.\ $X_k$ with mean $\mu$, let $Y_k=X_k\mathbf 1_{|X_k|\le k}$. By~\eqref{eq:etemadisum} (applied to $|X_k|$), $\sum_k\operatorname{Var}(Y_k/k)<\infty$, so
$\sum(Y_k-\mathbb E Y_k)/k$ converges almost surely, and Kronecker's lemma gives $\frac1n\sum_{k\le n}(Y_k-\mathbb E Y_k)\to0$. As in Etemadi's proof,
$\mathbb E Y_k\to\mu$ and $X_k=Y_k$ eventually. This proof uses full independence (through Proposition~\ref{prop:l2series}), but it gives more: by the
same route, $S_n/n^{1/p}\to0$ almost surely for centred $X$ with $\mathbb E|X|^p<\infty$, $1\le p<2$ (Exercise~\ref{ex:mz}).
\end{remark}

\section{Renewal}

A light bulb is replaced as soon as it fails; the lifetimes $\xi_1,\xi_2,\dots$ are i.i.d., positive, with mean $\mu\in(0,\infty]$. Let
$T_n=\xi_1+\dots+\xi_n$ be the time of the $n$-th replacement and $N(t)=\max\{n:T_n\le t\}$ the number of replacements by time $t$.

\begin{theorem}[strong law of renewal]\label{thm:renewal}
$N(t)/t\to1/\mu$ almost surely as $t\to\infty$ (with $1/\infty=0$).
\end{theorem}
\begin{proof}
By the strong law (Exercise~\ref{ex:infmeanpos} if $\mu=\infty$), $T_n/n\to\mu$ almost surely. Each $T_n$ is finite, so $N(t)\to\infty$ as
$t\to\infty$. By definition $T_{N(t)}\le t<T_{N(t)+1}$, hence
\[
\frac{T_{N(t)}}{N(t)}\le\frac t{N(t)}<\frac{T_{N(t)+1}}{N(t)+1}\cdot\frac{N(t)+1}{N(t)}.
\]
Both bounds tend to $\mu$ along the path, so $t/N(t)\to\mu$.
\end{proof}

The expectation $\mathbb E N(t)/t\to1/\mu$ also holds (the \emph{elementary renewal theorem}), but almost sure convergence does not give it for free;
Exercise~\ref{ex:wald} proves it with Wald's identity. Renewal processes count failures, arrivals at a server, and the cache misses of an algorithm,
and Theorem~\ref{thm:renewal} turns statements about one cycle into statements about long-run rates: the \emph{renewal--reward theorem} says that the long-run reward per unit time is the expected reward per cycle divided by the expected cycle length.

\begin{beauty}
The strong law is a statement about a single infinite experiment, and a theorem about averages over all of them; the bridge is measure theory.
Etemadi's proof shows how little it needs: variances of pairs, a sparse subsequence of times, and the monotonicity of sums of positive numbers. The
St Petersburg game shows the other side: a game with infinite expectation still has a law of large numbers, at the scale $n\log n$, but only in
probability. The difference between ``in probability'' and ``almost surely'' is the difference between what a casino can rely on and what a single
gambler will see.
\end{beauty}

\section{Exercises}

\begin{exercise}[$\star$]\label{ex:infmeanpos}
Let $X_k\ge0$ be i.i.d.\ with $\mathbb E X_1=\infty$. Show that $S_n/n\to\infty$ almost surely. (Truncate at a fixed level $M$.)
\end{exercise}

\begin{exercise}[$\star$]\label{ex:wald}
In the setting of Theorem~\ref{thm:renewal}, with $\mu<\infty$: (a) show that $N(t)+1$ is a stopping time for the filtration of the $\xi_k$, that
$\mathbb E N(t)<\infty$, and that $\mathbb E T_{N(t)+1}=\mu\,\mathbb E[N(t)+1]$ (Wald's identity; use the martingale $T_n-n\mu$ and optional stopping at
$(N(t)+1)\wedge m$). (b) Deduce $\liminf\mathbb E N(t)/t\ge1/\mu$ from $T_{N(t)+1}>t$. (c) Show $\limsup\mathbb E N(t)/t\le1/\mu$, first when $\xi_k\le M$,
then in general by truncating the lifetimes, which can only increase $N(t)$.
\end{exercise}

\begin{exercise}[$\star\star$]\label{ex:feller}
Let $X_k$ be i.i.d. (a) If $xP(|X|>x)\to0$ as $x\to\infty$, show that $S_n/n-\mathbb E[X\mathbf 1_{|X|\le n}]\to0$ in probability, using
$\mathbb E[X^2\mathbf 1_{|X|\le n}]=\int_0^n2yP(|X|>y)\dd y-n^2P(|X|>n)\le\int_0^n2yP(|X|>y)\dd y$. (b) Let $X$ be symmetric with $P(|X|>x)=\frac e{x\log x}$ for
$x\ge e$. Show that $S_n/n\to0$ in probability, but that $S_n/n$ does not converge almost surely.
\end{exercise}

\begin{exercise}[$\star\star$]\label{ex:stpeteas}
In the St Petersburg game, show that $\limsup_nS_n/(n\log_2n)=\infty$ almost surely. So Theorem~\ref{thm:stpetersburg} cannot be improved to almost sure
convergence.
\end{exercise}

\begin{exercise}[$\star\star$]\label{ex:seriesiid}
Let $X_k$ be i.i.d.\ and centred. (a) Show with the three-series theorem (with $c=1$, applied to $X_k/k$) that $\sum_kX_k/k$ converges almost surely
if $\mathbb E[|X_1|\log^+|X_1|]<\infty$. (b) Let $X=Y-\mathbb E Y$ where $P(Y>y)=\frac e{y(\log y)^2}$ for $y\ge e$. Show that $\mathbb E|X|<\infty$ but that
$\sum_kX_k/k$ diverges almost surely.
\end{exercise}

\begin{exercise}[$\star\star\star$]\label{ex:mz}
Let $1\le p<2$ and let $X_k$ be i.i.d., centred, with $\mathbb E|X_1|^p<\infty$. Show that $S_n/n^{1/p}\to0$ almost surely (Marcinkiewicz and
Zygmund, 1937). (Apply the three-series theorem to $X_k/k^{1/p}$, then Kronecker's lemma.)
\end{exercise}

\begin{exercise}[$\star\star\star$]\label{ex:chungfuchs}
Let $S_n$ be a random walk on $\Z$ with i.i.d.\ steps, $\mathbb E|X_1|<\infty$ and $\mathbb E X_1=0$. Let $G_N(x)=\sum_{n=0}^NP(S_n=x)$. (a) Show that
$G_N(x)\le G_N(0)$ for every $x$ (decompose according to the first visit to $x$). (b) Using the weak law, show that
$\sum_{|x|\le\varepsilon N}G_N(x)\ge\sum_{n\le N}P(|S_n|\le\varepsilon n)\ge(1-o(1))N$. (c) Deduce that $\sum_nP(S_n=0)=\infty$, so that the walk is recurrent
(Tome~3, Theorem~7.12). This is the theorem of Chung and Fuchs (1951) in dimension one: every centred integrable walk on $\Z$ returns to $0$.
\end{exercise}

\section{Hints}

\begin{hint}{ex:infmeanpos}
For fixed $M$, $S_n/n$ is at least the average of the $\min(X_k,M)$, $k\le n$, which tends to $\mathbb E\min(X_1,M)$, which tends to $\infty$ with $M$.
\end{hint}
\begin{hint}{ex:wald}
$\{N(t)+1=n\}=\{T_{n-1}\le t<T_n\}$. For (c) with $\xi\le M$: $T_{N(t)+1}\le t+M$.
\end{hint}
\begin{hint}{ex:feller}
(a) Theorem~\ref{thm:wlltriangular} with $b_n=n$; the average of a function tending to $0$ tends to $0$. (b) The mean of the truncation is $0$ by symmetry;
for almost sure convergence, use Exercise~\ref{ex:infmean} of Chapter~1.
\end{hint}
\begin{hint}{ex:stpeteas}
$P(X_n>Cn\log_2n)\ge\frac1{Cn\log_2n}$, and $\sum\frac1{n\log n}=\infty$.
\end{hint}
\begin{hint}{ex:seriesiid}
$\sum_kP(|X_1|>k)<\infty$; for (iii) use~\eqref{eq:etemadisum}; for (ii), $\mathbb E[X\mathbf 1_{|X|\le k}]=-\mathbb E[X\mathbf 1_{|X|>k}]$. In (b), series (ii)
is $-\sum_k\frac1k\mathbb E[X\mathbf 1_{|X|>k}]$ and $\mathbb E[X\mathbf 1_{X>k}]\approx\frac e{\log k}$.
\end{hint}
\begin{hint}{ex:mz}
The three series are $\sum P(|X|^p>k)$, $\sum k^{-2/p}\mathbb E[X^2\mathbf 1_{|X|\le k^{1/p}}]$ and $\sum k^{-1/p}\mathbb E[X\mathbf 1_{|X|>k^{1/p}}]$ (up to sign).
\end{hint}
\begin{hint}{ex:chungfuchs}
(a) Let $\tau_x$ be the first visit to $x$; after it the walk starts afresh from $x$. (c) Combine (a) and (b): $(2\varepsilon N+1)G_N(0)\ge(1-o(1))N$.
\end{hint}

\section{Solutions}

\begin{solution}{ex:infmeanpos}
\idea Infinite mean means the truncated means are unbounded; the strong law applies to each truncation, and the truncation only lowers the sum.
\writeup For fixed $M$, the variables $\min(X_k,M)$ are i.i.d.\ and bounded, so by Theorem~\ref{thm:etemadi}, $\frac1n\sum_{k\le n}\min(X_k,M)\to\mathbb E\min(X_1,M)$
almost surely. Since $X_k\ge\min(X_k,M)$, $\liminf S_n/n\ge\mathbb E\min(X_1,M)$ almost surely, for every $M$ (countably many $M=1,2,\dots$ suffice). By monotone
convergence $\mathbb E\min(X_1,M)\to\mathbb E X_1=\infty$.
\end{solution}

\begin{solution}{ex:wald}
\idea Wald's identity is optional stopping for the martingale $T_n-n\mu$; the stopping time $N(t)+1$ is the first $n$ with $T_n>t$, and the overshoot
$T_{N(t)+1}-t$ is at most one lifetime.
\writeup (a) $\{N(t)+1\le n\}=\{T_n>t\}$ depends on $\xi_1,\dots,\xi_n$, so $\sigma=N(t)+1$ is a stopping time. Choose $\delta>0$ with $p=P(\xi>\delta)>0$; the
number of lifetimes exceeding $\delta$ before time $t$ is at most $t/\delta$, and $N(t)$ is at most the number of trials needed for $\lfloor t/\delta\rfloor+1$
successes of probability $p$, a negative binomial variable with finite mean. So $\mathbb E N(t)<\infty$. Optional stopping at the bounded time $\sigma\wedge m$ gives
$\mathbb E T_{\sigma\wedge m}=\mu\,\mathbb E[\sigma\wedge m]$, and both sides increase to their limits by monotone convergence, so $\mathbb E T_\sigma=\mu\,\mathbb E\sigma$.

(b) $T_\sigma>t$, so $\mu(\mathbb E N(t)+1)>t$, i.e.\ $\mathbb E N(t)/t>1/\mu-1/t$.

(c) If $\xi_k\le M$, then $T_\sigma\le t+M$, so $\mu(\mathbb E N(t)+1)\le t+M$ and $\limsup\mathbb E N(t)/t\le1/\mu$. In general, let $\xi_k^M=\min(\xi_k,M)$ with mean
$\mu_M$ and renewal count $N^M(t)\ge N(t)$ (shorter lifetimes, more replacements). Then $\limsup\mathbb E N(t)/t\le1/\mu_M$, and $\mu_M\uparrow\mu$ as $M\to\infty$.
\end{solution}

\begin{solution}{ex:feller}
\idea Feller's condition $xP(|X|>x)\to0$ is exactly what makes both conditions of the triangular weak law work at the scale $b_n=n$. The example
satisfies it while having an infinite mean, so the strong law must fail.
\writeup (a) Condition (a) of Theorem~\ref{thm:wlltriangular} is $nP(|X|>n)\to0$. For (b), the formula $\mathbb E[Z^2]=\int_0^\infty2yP(|Z|>y)\dd y$ applied to
$Z=X\mathbf 1_{|X|\le n}$ (for which $P(|Z|>y)\le P(|X|>y)$ and $=0$ for $y\ge n$) gives $\frac1n\mathbb E[X^2\mathbf 1_{|X|\le n}]\le\frac1n\int_0^n2yP(|X|>y)\dd y$.
The integrand $g(y)=2yP(|X|>y)$ tends to $0$ and is bounded by $2y$, so its average over $[0,n]$ tends to $0$. Hence $(S_n-n\,\mathbb E[X\mathbf 1_{|X|\le n}])/n\to0$
in probability.

(b) At $x=e$, $P(|X|>e)=1$, so this is a valid law. Here $xP(|X|>x)=e/\log x\to0$ and the truncated means vanish by symmetry, so $S_n/n\to0$ in probability.
But $\mathbb E|X|=\int_0^\infty P(|X|>x)\dd x\ge\int_e^\infty\frac{e\dd x}{x\log x}=\infty$, so by Exercise~\ref{ex:infmean} of Chapter~1, $\limsup|S_n|/n=\infty$
almost surely: $S_n/n$ converges nowhere, except on a null set.
\end{solution}

\begin{solution}{ex:stpeteas}
\idea A single prize larger than $Cn\log_2n$ at time $n$ already pushes the ratio above $C$; such prizes happen infinitely often because the
probabilities form a divergent series.
\writeup Since $P(X\ge2^j)=2^{1-j}$, for $x\ge2$ we have $P(X>x)\ge P(X\ge2^{\lceil\log_2x\rceil+1})\ge\frac1{2x}$. Hence
$\sum_nP(X_n>Cn\log_2n)\ge\sum_n\frac1{2Cn\log_2n}=\infty$, and the events are independent; by Borel--Cantelli, $X_n>Cn\log_2n$ for infinitely many $n$. For each
such $n$, $S_n\ge X_n>Cn\log_2n$. So $\limsup S_n/(n\log_2n)\ge C$ almost surely, for every $C$.
\end{solution}

\begin{solution}{ex:seriesiid}
\idea The three series for $X_k/k$ with $c=1$ involve the tails of $X$ at the levels $k$, exactly the quantities controlled in Etemadi's proof. The
delicate one is the series of means: the truncation at $k$ removes a little mass from the tail, and the removed amounts, divided by $k$, must sum.
\writeup (a) (i) $\sum_kP(|X_k/k|>1)=\sum_kP(|X_1|>k)\le\mathbb E|X_1|<\infty$. (iii) The variance of $X_k\mathbf 1_{|X_k|\le k}/k$ is at most $k^{-2}\mathbb E[X_1^2\mathbf 1_{|X_1|\le k}]$,
and the sum over $k$ is at most $2\mathbb E|X_1|$ by~\eqref{eq:etemadisum} applied to $|X_1|$. (ii) Since $\mathbb E X_1=0$,
\[
\bigl|\mathbb E[X_1\mathbf 1_{|X_1|\le k}]\bigr|=\bigl|\mathbb E[X_1\mathbf 1_{|X_1|>k}]\bigr|\le\mathbb E[|X_1|\mathbf 1_{|X_1|>k}],
\]
so
\[
\begin{gathered}
\sum_k\frac1k\bigl|\mathbb E[X_1\mathbf 1_{|X_1|\le k}]\bigr|\le\mathbb E\Bigl[|X_1|\sum_{k<|X_1|}\frac1k\Bigr]\qq\le\mathbb E\bigl[|X_1|(1+\log^+|X_1|)\bigr]<\infty.
\end{gathered}
\]
All three series converge.

(b) $\mathbb E Y=e+\int_e^\infty\frac{e\dd y}{y(\log y)^2}=2e<\infty$, so $X$ is integrable and centred. For $k\ge4e$, $|X_1|>k$ means $Y_1>k+2e$ (since $X_1\ge e-2e>-k$), and on that event
$Y_1-2e\ge\frac12Y_1$; using $\mathbb E[Y\mathbf 1_{Y>t}]\ge\int_t^\infty P(Y>y)\dd y$,
\[
\mathbb E[X_1\mathbf 1_{|X_1|>k}]=\mathbb E[(Y_1-2e)\mathbf 1_{Y_1>k+2e}]\ge\tfrac12\mathbb E[Y_1\mathbf 1_{Y_1>k+2e}]\ge\tfrac12\int_{k+2e}^\infty\frac{e\dd y}{y(\log y)^2}
=\frac e{2\log(k+2e)}.
\]
So the series (ii), which is
\[
\sum_k\frac1k\mathbb E[X_1\mathbf 1_{|X_1|\le k}]=-\sum_k\frac1k\mathbb E[X_1\mathbf 1_{|X_1|>k}],
\]
diverges like $-\sum\frac e{2k\log k}$. By the
three-series theorem, $\sum X_k/k$ diverges almost surely. Here $\mathbb E[|X|\log^+|X|]=\infty$: the condition of (a) is sharp for this kind of one-sided tail.
\end{solution}

\begin{solution}{ex:mz}
\idea Rescale: $X_k/k^{1/p}$ is the right summand, and each of the three series turns into a moment of $|X|^p$ by exchanging sum and expectation, as
in~\eqref{eq:etemadisum}.
\writeup Let $c=1$, $Z_k=X_k/k^{1/p}$, $W_k=Z_k\mathbf 1_{|Z_k|\le1}$. (i) $\sum_kP(|X_1|>k^{1/p})=\sum_kP(|X_1|^p>k)\le\mathbb E|X_1|^p$.

(iii) $\sum_k\mathbb E W_k^2=\mathbb E\bigl[X_1^2\sum_{k\ge|X_1|^p}k^{-2/p}\bigr]$. For $y\ge1$, $\sum_{k\ge y}k^{-2/p}\le y^{-2/p}+\int_y^\infty x^{-2/p}\dd x=
y^{-2/p}+\frac{y^{1-2/p}}{2/p-1}\le C_py^{1-2/p}$, since $2/p>1$. With $y=\max(|X_1|^p,1)$ the bracket is at most $C_p|X_1|^2\max(|X_1|^p,1)^{1-2/p}\le C_p\max(|X_1|^p,|X_1|^2)$,
and when $|X_1|<1$ this is at most $C_p$; in all cases at most $C_p(1+|X_1|^p)$, which is integrable.

(ii) Since $\mathbb E X_1=0$,
\[
|\mathbb E W_k|=k^{-1/p}\bigl|\mathbb E[X_1\mathbf 1_{|X_1|>k^{1/p}}]\bigr|\le k^{-1/p}\mathbb E[|X_1|\mathbf 1_{|X_1|^p>k}],
\]
so $\sum_k|\mathbb E W_k|\le\mathbb E\bigl[|X_1|\sum_{k<|X_1|^p}k^{-1/p}\bigr]$. When $p>1$, $\sum_{k<y}k^{-1/p}\le\frac{y^{1-1/p}}{1-1/p}$, and the bracket is at most
$\frac p{p-1}|X_1|\cdot|X_1|^{p-1}=\frac p{p-1}|X_1|^p$: integrable. The case $p=1$ is Theorem~\ref{thm:etemadi} itself.

So $\sum_kX_k/k^{1/p}$ converges almost surely for $1<p<2$, and Kronecker's lemma with $b_k=k^{1/p}$ gives $S_n/n^{1/p}\to0$.
\end{solution}

\begin{solution}{ex:chungfuchs}
\idea Recurrence on $\Z$ is $\sum_nP(S_n=0)=\infty$. The weak law says the walk spends most of its first $N$ steps in the window $|x|\le\varepsilon N$, of
about $2\varepsilon N$ sites; if no site is visited more often than $0$, then $0$ must be visited at least about $\frac1{2\varepsilon}$ times.
\writeup (a) Let $\tau_x=\min\{n\ge0:S_n=x\}$. The event $\{\tau_x=m\}$ depends only on $X_1,\dots,X_m$, so on it the walk $(S_{m+j}-x)_{j\ge0}=(X_{m+1}+\dots+X_{m+j})_{j\ge0}$ is a
fresh copy of the walk, independent of $\{\tau_x=m\}$ (the strong Markov property at $\tau_x$). So
\[
G_N(x)=\sum_{m\le N}P(\tau_x=m)\sum_{j\le N-m}P(S_j=0)\le\sum_mP(\tau_x=m)\,G_N(0)\le G_N(0).
\]
(b) $\sum_{|x|\le\varepsilon N}G_N(x)=\sum_{n\le N}P(|S_n|\le\varepsilon N)\ge\sum_{n\le N}P(|S_n|\le\varepsilon n)$. By the weak law, $P(|S_n|\le\varepsilon n)\to1$,
so the Ces\`aro average $\frac1N\sum_{n\le N}P(|S_n|\le\varepsilon n)\to1$.

(c) By (a) the left side of (b) is at most $(2\varepsilon N+1)G_N(0)$. Hence $G_N(0)\ge\frac{(1-o(1))N}{2\varepsilon N+1}$, and $\liminf_NG_N(0)\ge\frac1{2\varepsilon}$.
As $\varepsilon$ is arbitrary, $\sum_nP(S_n=0)=\infty$, and Tome~3, Theorem~7.12 says the walk returns to $0$ infinitely often, almost surely.
\end{solution}

\section*{Sources and further reading}

\begin{itemize}
\item \emph{History.} J.~Bernoulli, \emph{Ars Conjectandi}, Basel, 1713, Part~IV; the computation of $25\,550$ draws and the later history are in E.~Seneta,
``A tricentenary history of the law of large numbers'', \emph{Bernoulli} 19 (2013), 1088--1121, which also records Poisson's name for the law (\emph{Recherches
sur la probabilit\'e des jugements}, 1837). F.~P.~Cantelli, ``Sulla probabilit\`a come limite della frequenza'', \emph{Rend.\ Accad.\ Lincei} 26 (1917), 39--45.
A.~Khintchine, ``Sur la loi des grands nombres'', \emph{C.\ R.\ Acad.\ Sci.\ Paris} 188 (1929), 477--479. A.~N.~Kolmogorov, ``Sur la loi forte des grands nombres'',
\emph{C.\ R.\ Acad.\ Sci.\ Paris} 191 (1930), 910--912. N.~Etemadi, ``An elementary proof of the strong law of large numbers'', \emph{Z.\ Wahrsch.\ Verw.\ Gebiete}
55 (1981), 119--122. The St Petersburg game: D.~Bernoulli, ``Specimen theoriae novae de mensura sortis'', \emph{Comment.\ Acad.\ Sci.\ Imp.\ Petropolitanae} 5 (1738),
175--192, which quotes Nicolaus Bernoulli's problem; W.~Feller, ``Note on the law of large numbers and `fair' games'', \emph{Ann.\ Math.\ Statist.}\ 16 (1945), 301--304;
Y.~S.~Chow and H.~Robbins, ``On sums of independent random variables with infinite moments and `fair' games'', \emph{Proc.\ Nat.\ Acad.\ Sci.\ USA} 47 (1961), 330--335.
\item R.~Durrett, \emph{Probability: Theory and Examples}, 5th ed., Cambridge University Press, 2019, Sections~2.2--2.5 (the triangular weak law, St Petersburg,
Etemadi, three series, renewal) and~5.4 (Chung--Fuchs).
\item W.~Feller, \emph{An Introduction to Probability Theory and Its Applications}, vol.~II, 2nd ed., Wiley, 1971, Chapters~VII and~XI. Laws of large numbers and
renewal theory.
\item J.~Marcinkiewicz and A.~Zygmund, ``Sur les fonctions ind\'ependantes'', \emph{Fund.\ Math.}\ 29 (1937), 60--90.
\item K.~L.~Chung and W.~H.~J.~Fuchs, ``On the distribution of values of sums of random variables'', \emph{Mem.\ Amer.\ Math.\ Soc.}\ 6 (1951).
\end{itemize}
